\section{Fixed-source closure and the finite-section edge theorem}\label{sec:edge}
\subsection{An invariant density class}
Choose $R$ with $2h<R<5h/2$. Let $\mathcal C$ consist of functions
$g$ whose density $\rho_g=w_yg$ is holomorphic on $|\Im y|<R$ and,
on every closed narrower strip, satisfies for some integer $M\ge0$
\begin{align}
 \rho_g(y)&=c_ge^{-ay}+O((1+\Re y)^Me^{-2a\Re y})
                                  &&(\Re y\to+\infty),\label{densityclass}\\
 \rho_g(y)&=O((1+|\Re y|)^Me^{2a\Re y})
                                  &&(\Re y\to-\infty).\nonumber
\end{align}
Such functions lie in $\mathcal H$ by the same tail calculation
used for \eqref{qsource}. Cauchy's formula supplies the corresponding
fixed-derivative estimates on narrower strips.

Both $q_3$ and $q^{(1)}$ belong to this class. In
\eqref{rhosource} the affine numerator vanishes identically in $t$
at $y=\pm ih/2,\pm3ih/2$, as direct substitution shows. These
remove the first four zeros of $2\cosh(3ay)$; the next possible
poles are $\pm5ih/2$. The tails give
$c_{q_3}=\sqrt{3/2}$ and $c_{q^{(1)}}=c_{q_3}/4$.

The exact adjoint density formula is
\begin{equation}\label{adjointdensity}
 \rho_{\mathcal B^*g}(y)
   =e^{ay}\int_\R\varphi(y-x)e^{-ax}m(x)\rho_g(x)\,dx.
\end{equation}
It follows directly from the unweighted kernel, first on the real
axis by absolute convergence. Write $F(x)=e^{-ax}m(x)\rho_g(x)$.
It is holomorphic on $|\Im x|<3h/2$, the nearest poles of $m$,
and its horizontal-strip tails are
\[
 F(x)=O(\operatorname{poly}(x)e^{-4a\Re x})\quad(+\infty),\qquad
 F(x)=O(\operatorname{poly}(|x|)e^{a\Re x})\quad(-\infty).
\]
The kernel is holomorphic for $|\Im u|<h$ and has tails
$O(e^{-2a|\Re u|})$ on narrower strips. To extend
\eqref{adjointdensity} to $|\Im y|<5h/2$, choose a contour height
$v$ with $|v|<3h/2$ and $|\Im y-v|<h$. Overlaps agree by contour
shifting; the vertical sides vanish by the displayed tails.
A finite cover gives uniform estimates on each narrower closed strip.

The positive leading coefficient is exactly
\begin{equation}\label{leadingmoment}
 c_{\mathcal B^*g}=\tfrac13\int_\R e^{ax}m(x)\rho_g(x)\,dx.
\end{equation}
To see the error rigorously, subtract $e^{-2a(y-x)}/3$ from the
kernel over the entire chosen contour. For $\Re x\le\Re y$ its
remainder is $O(e^{-4a(\Re y-\Re x)})$; for $\Re x>\Re y$
bound the kernel and subtracted exponential separately. Combining
with the tails of $F$ gives a polynomial times $e^{-4a\Re y}$
before multiplication by $e^{ay}$. The borderline integration over
$0<\Re x<\Re y$ only adds one polynomial power. The final remainder
is $O(\operatorname{poly}(y)e^{-3a\Re y})$, stronger than needed.
The moment integrand in \eqref{leadingmoment} is analytic in the
contour strip and decays at both ends, so its value is contour
independent. On the negative tail split at $\Re x=\Re y$ and
$\Re x=0$. The convolution is
$O(\operatorname{poly}(|y|)e^{a\Re y})$ before multiplication,
giving the required $e^{2a\Re y}$ bound. Thus
\begin{equation}\label{classclosure}
 \mathcal B^*\mathcal C\subset\mathcal C.
\end{equation}

Every member has the same leading coefficient shape:
\begin{equation}\label{shape}
 \langle\psi_j,g\rangle
   =\frac{c_g}{c_{q_3}}a_j^{\rm src}+o(1/(j+1)).
\end{equation}
Here is a proof of the little-o claim, including the endpoint domain.
Subtract $(c_g/c_{q_3})q_3$ and call the result $d_g$. Its density
has two-sided $\operatorname{poly}(y)e^{-2a|\Re y|}$ tails in the
strip. In $z=y/(3h)$ put
$G^\sharp(z)=\Gamma(5/6-iz)\Gamma(7/6+iz)$.
The exactly normalized Fourier input in the Jacobi transform is
$\sqrt{c_W}G(z)d_g(3hz)$, $c_W=81\sqrt3/(8\pi^2)$.
Since $w_y(3hz)$ is a constant times $G(z)G^\sharp(z)$, this input
is a fixed constant times $\rho_{d_g}(3hz)/G^\sharp(z)$.
The reciprocal gamma product is entire. Stirling on horizontal strips
makes this quotient, and all its fixed polynomial multiples, exponentially
decaying. Its holomorphy strip contains $\Im z=\pm2/3$, because
$R>2h$. Contour shifting for its Fourier transform $H_g$ yields
\[
 |H_g(\xi)|+|H_g'(\xi)|\le Ce^{-2|\xi|/3}.
\]
The angular image
$g_\theta(\theta)=\sqrt{2/\sin\theta}\,H_g(\xi(\theta))$
therefore satisfies
$g_\theta=O(d^{5/6})$ and $g_\theta'=O(d^{-1/6})$.
It belongs to $H_0^1(0,\pi)$ and
$\int |g_\theta|^2/d^2<\infty$. Hence it is in the quadratic-form
domain of the Friedrichs Jacobi angular operator, with potential
\[
 \frac{\alpha_J^2-1/4}{4\sin^2(\theta/2)}
 +\frac{\beta_J^2-1/4}{4\cos^2(\theta/2)},
 \qquad (\alpha_J,\beta_J)=(4/3,2/3).
\]
This is the realization with eigenfunctions $u_j$ and eigenvalues
$(j+3/2)^2$, also at the limit-circle endpoint.
Its form expansion gives
$\sum_j(j+3/2)^2|\langle\psi_j,d_g\rangle|^2<\infty$;
the unitary changes coefficients only by phases $i^j$.
Thus $j\langle\psi_j,d_g\rangle\to0$, proving \eqref{shape}.
In particular these coefficients are $O((j+1)^{-1})$ and their
projection-tail norms are $O(n^{-1/2})$.
By \eqref{classclosure}, all fixed iterates of both $q_3$ and
$q^{(1)}$ have this shape. The constants in \eqref{leadingmoment}
retain the full operator, including its compact and rank-one pieces.

\subsection{An abstract edge theorem with nonlocal tails retained}
\begin{theorem}\label{thm:edge}
Let $B$ be bounded on $\ell^2(\N_0)$ and $T_b=T_b^*$ a Toeplitz
operator with entries $b_{i-k}$, where for some $\eps>0$
\[
 |b_r|\le C(1+|r|)^{-1-\eps},\qquad
 \|(B-T_b)e_i\|=o((i+1)^{-1/2}).
\]
Let $q$ be fixed and $u_\ell=(B^*)^\ell q$. Suppose for each fixed
$\ell\ge0$ that
\[
 (u_\ell)_i=\kappa_\ell a_i+o((i+1)^{-1}),\qquad
 a_i\sim c_*( -1)^i/(i+1),\quad c_*\ne0.
\]
Then for every fixed $\ell\ge0$ the full-sequence limit
\[
 \frac{[(P_nB^*P_n)^\ell P_nq]_{n-1}}{a_{n-1}}
\]
exists. Every fixed right-edge profile coordinate has a limit as well.
\end{theorem}
\begin{proof}
Put $Q_n=I-P_n$, $A_n=P_nB^*P_n$, and
$e_{n,\ell}=A_n^\ell P_nq-P_nu_\ell$. Exactly
\begin{equation}\label{defectrecurrence}
 e_{n,0}=0,\qquad
 e_{n,\ell+1}=P_nB^*e_{n,\ell}-P_nB^*Q_nu_\ell.
\end{equation}
The coefficient hypothesis gives $\|Q_nu_\ell\|=O(n^{-1/2})$,
and boundedness of $B$ then gives
$\|e_{n,\ell}\|=O(n^{-1/2})$ for each fixed $\ell$.
Set $K=B-T_b$. If $\|v_n\|=O(n^{-1/2})$, the column hypothesis
implies, uniformly for $i\ge\eta n$, $\eta>0$ fixed,
\begin{equation}\label{compacttail}
 |(K^*v_n)_i|\le\|Ke_i\|\|v_n\|=o(1/n).
\end{equation}
Uniformity follows because
$\sup_{i\ge\eta n}\sqrt{i+1}\|Ke_i\|\to0$.
This uses columns of $K$ to control coordinates of $K^*$, without
assuming an adjoint-column estimate.

Induction in \eqref{defectrecurrence} gives the stronger pointwise bound
\begin{equation}\label{pointwisedefect}
 \sup_{\eta n\le i<n}n|(e_{n,\ell})_i|\le C_{\ell,\eta}.
\end{equation}
For its Toeplitz part split the input at $\eta n/2$. On the upper
part use the preceding inductive bound and $\sum|b_r|<\infty$.
On the lower part the distance is at least $\eta n/2$;
Cauchy--Schwarz bounds the Toeplitz row tail by
$O(n^{-1/2-\eps})$, which times $\|e_{n,\ell}\|=O(n^{-1/2})$
is $O(n^{-1-\eps})$. The source tail $Q_nu_\ell$ has coordinates
$O(1/n)$ for all $i\ge n$, so its Toeplitz image is $O(1/n)$.
The two $K^*$ terms satisfy \eqref{compacttail}. This proves
\eqref{pointwisedefect}, starting from zero.

For each fixed $r\ge0$ let
$E_\ell(r)=\lim_{n\to\infty}(e_{n,\ell})_{n-1-r}/a_{n-1}$
when it exists. It is zero for $\ell=0$. Assuming it exists at
level $\ell$, the exact recurrence gives
\begin{equation}\label{edgeprofile}
 E_{\ell+1}(r)=\sum_{s\ge0}b_{s-r}E_\ell(s)
       -\kappa_\ell\sum_{m\ge0}b_{-r-1-m}(-1)^{m+1}.
\end{equation}
To justify all tails, on input indices $k\ge n/2$ the normalized
defect coordinates are uniformly bounded by \eqref{pointwisedefect};
absolute summability permits a fixed-distance cutoff followed by its
removal. On $k<n/2$, the preceding row-tail estimate is
$O(n^{-1-\eps})=o(a_{n-1})$. On the exterior tail $k=n+m$, the
normalized source coordinates are uniformly bounded and converge,
for each fixed $m$, to $\kappa_\ell(-1)^{m+1}$. Dominated summation
therefore applies to the second series. The $K^*$ terms are
$o(a_{n-1})$ by \eqref{compacttail}. Each limiting profile is bounded
by \eqref{pointwisedefect} with $\eta=1/2$.
This proves existence inductively and gives
\[
 \lim_{n\to\infty}\frac{[A_n^\ell P_nq]_{n-1}}{a_{n-1}}
       =\kappa_\ell+E_\ell(0).
\]
There is no infinite-power interchange in this argument.
\end{proof}

\subsection{Closing the analytic boundary limit}
Apply Theorem~\ref{thm:edge} to the original $B$ in the $\psi$ basis,
using Theorem~\ref{thm:columns} with $\eps=2/3$ and the source shape
\eqref{shape}. Apply it separately to $q_3$ and $q^{(1)}$.
Every coefficient in \eqref{Taylorcoef} therefore converges along
the full sequence. The normal-family argument following that formula
proves $F_n\to F$ locally uniformly on $\HH_+$.
The companion identity gives exactly
\[
 F_n(s)F_{n+1}(s)
  =(s/c)^2\frac{D_n(s^2)/D_n(3)}{D_{n+1}(s^2)/D_{n+1}(3)}.
\]
Since $E_n\to E$ locally uniformly, the right side tends to
$(s/c)^2/P(s)$. Thus $F(s)^2=(s/c)^2/P(s)$.
It is nonzero on $\HH_+$, and $F(c)=1$ selects the branch
$F(s)=(s/c)/\sqrt{P(s)}$. This proves \eqref{Ftarget} without a
residual parity ambiguity.

\subsection{A nonzero leading Galerkin contribution as a check}
There is a useful independent check on the sign and scale.
Define
\[
 \tau_j=\frac{\langle B e_j,Q_{j+1}q_3\rangle}{a_j^{\rm src}},
\]
with source and basis now in coordinate space. The $B-T_b$ part is
$o(1)$ after division by $a_j^{\rm src}$, by
\eqref{originalcolumns} and $\|Q_{j+1}q_3\|=O(j^{-1/2})$.
For the Toeplitz part, the calibration asymptotic makes
$a_{j+r}^{\rm src}/a_j^{\rm src}$ uniformly bounded for $r\ge1$
and gives its fixed-$r$ limit $(-1)^r$. Absolute summability yields
\begin{equation}\label{taulimit}
 \tau_j\longrightarrow\sum_{r\ge1}(-1)^rb_r=\sqrt3/4-1/2.
\end{equation}
Thus the original finite-projection tail is genuinely leading order.
It was retained in \eqref{edgeprofile}; compactness alone would not
have justified discarding it.
